% !TEX root = ../main.tex
%% Chapitre 18 — Espaces préhilbertiens réels
%% Programme officiel MPSI/MP2I (BO spécial n°1 du 11/02/2021), deuxième semestre :
%% a) produit scalaire (canonique, matriciel tr(AᵀB), intégral) ; b) norme,
%% Cauchy-Schwarz et triangulaire avec cas d'égalité, polarisation ;
%% c) orthogonalité, familles orthonormées, Pythagore, Gram-Schmidt ;
%% d) bases orthonormées ; e) projection orthogonale sur un sous-espace de
%% dimension finie, distance, vecteur normal d'un hyperplan.
%% Résultats vérifiés par verifications/chapitre_18.py.

\begin{ouverture}[orthogonalite]{18}{Espaces préhilbertiens réels}{Partie 2 · Deuxième semestre}
  {Produit scalaire\enspace·\enspace Cauchy-Schwarz\enspace·\enspace Orthogonalité\enspace·\enspace
   Gram-Schmidt\enspace·\enspace Projection et distance}
  \objectif{Reconnaître un produit scalaire : sur $\R^{n}$, sur les matrices, sur les fonctions.}
  \objectif{Manier la norme : Cauchy-Schwarz et l'inégalité triangulaire, avec leurs cas d'égalité.}
  \objectif{Exploiter l'orthogonalité : Pythagore, familles orthonormées, orthogonal d'une partie.}
  \objectif{Orthonormaliser une famille par l'algorithme de Gram-Schmidt.}
  \objectif{Travailler en base orthonormée : coordonnées, produit scalaire, norme.}
  \objectif{Projeter orthogonalement et calculer des distances — la meilleure approximation.}
  \prerequis{Espaces vectoriels, bases, dimension, supplémentaires (chapitre 12).}
  \prerequis{Trace et transposée (chapitre 13) ; intégrale sur un segment (chapitre 15).}
  \prerequis{Géométrie du plan et de l'espace vues au lycée — qu'on généralise ici.}
\end{ouverture}

%% =====================================================================
%%  L'ESSENTIEL
%% =====================================================================
\essentiel

\essSection{1}{Le produit scalaire}

\begin{definition}[Produit scalaire ; espaces préhilbertiens]
Un \textbf{produit scalaire} sur un $\R$-espace vectoriel $E$ est une
application $(x,y)\mapsto\langle x,y\rangle$ \emph{bilinéaire},
\emph{symétrique}, \emph{définie positive} :
$\langle x,x\rangle\geqslant0$, avec égalité seulement pour $x=0$. On note
aussi $(x\mid y)$ ou $x\cdot y$. Muni d'un produit scalaire, $E$ est dit
\textbf{préhilbertien} ; \textbf{euclidien} s'il est de dimension finie.
\end{definition}

\begin{exemple}[Les trois produits scalaires du programme]
Sur $\R^{n}$ : $\langle X,Y\rangle=\sum x_iy_i=X^{\mathsf T}Y$ (canonique).
Sur $\mathcal{M}_{n,p}(\R)$ :
$\langle A,B\rangle=\operatorname{tr}\bigl(A^{\mathsf T}B\bigr)
=\sum_{i,j}a_{ij}b_{ij}$ — déjà rencontré au chapitre 13. Sur
$\mathcal{C}\bigl(\intFF{a}{b},\R\bigr)$ ou sur $\R[X]$ :
$\langle f,g\rangle=\displaystyle\int_{a}^{b}fg$ — le caractère défini
positif est exactement l'exercice « intégrale nulle » du chapitre 15.
Trois mondes, une seule géométrie.
\end{exemple}

\essSection{2}{Norme, Cauchy-Schwarz, triangulaire}

\begin{definition}[Norme associée]
$\norme{x}=\sqrt{\langle x,x\rangle}$ ; la \emph{distance} associée est
$d(x,y)=\norme{x-y}$. Identité remarquable :
\[
\norme{x+y}^{2}=\norme{x}^{2}+\norme{y}^{2}+2\langle x,y\rangle ,
\]
d'où la \emph{polarisation}
$\langle x,y\rangle=\frac14\bigl(\norme{x+y}^{2}-\norme{x-y}^{2}\bigr)$ :
le produit scalaire se lit sur la norme — et l'identité du
\emph{parallélogramme}
$\norme{x+y}^{2}+\norme{x-y}^{2}=2\norme{x}^{2}+2\norme{y}^{2}$.
\end{definition}

\begin{theoreme}[Cauchy-Schwarz ; inégalité triangulaire]
Pour tous $x,y\in E$ :
\[
|\langle x,y\rangle|\leqslant\norme{x}\,\norme{y}
\qquad\text{et}\qquad
\norme{x+y}\leqslant\norme{x}+\norme{y} ,
\]
avec égalité dans la première si et seulement si $(x,y)$ est \emph{liée},
dans la seconde si et seulement si $x$ et $y$ sont colinéaires \emph{de
même sens}. La preuve de Cauchy-Schwarz est le trinôme
$t\mapsto\norme{x+ty}^{2}\geqslant0$, déjà mené pour les matrices
(chapitre 13) et les intégrales (chapitre 15) : c'était chaque fois le
même théorème.
\end{theoreme}

\begin{exemple}
Dans $\R^{n}$ avec $y=(1,\dots,1)$ :
$\bigl(\sum a_i\bigr)^{2}\leqslant n\sum a_i^{2}$ ; en intégral :
$\bigl(\int_0^1f\bigr)^{2}\leqslant\int_0^1f^{2}$ — la moyenne au carré ne
dépasse pas la moyenne des carrés.
\end{exemple}

\essSection{3}{Orthogonalité}

\begin{definition}[Vecteurs et familles orthogonaux]
$x\perp y$ si $\langle x,y\rangle=0$. L'\textbf{orthogonal} d'une partie
$A$ est $A^{\perp}=\ens{x\mid\forall a\in A,\ \langle x,a\rangle=0}$ :
c'est toujours un sous-espace vectoriel. Une famille est
\emph{orthogonale} si ses vecteurs sont deux à deux orthogonaux,
\emph{orthonormée} si de plus ils sont unitaires :
$\langle e_i,e_j\rangle=\delta_{ij}$.
\end{definition}

\begin{theoreme}[Liberté et Pythagore]
Toute famille orthogonale de vecteurs \emph{non nuls} est libre. Et si
$x_1,\dots,x_p$ sont deux à deux orthogonaux :
\[
\norme{x_1+\dots+x_p}^{2}=\norme{x_1}^{2}+\dots+\norme{x_p}^{2}
\qquad\text{(Pythagore).}
\]
Pour deux vecteurs, la réciproque est vraie :
$\norme{x+y}^{2}=\norme{x}^{2}+\norme{y}^{2}\iff x\perp y$.
\end{theoreme}

\essSection{4}{L'algorithme de Gram-Schmidt}

\begin{theoreme}[Orthonormalisation]
Soit $(v_1,\dots,v_p)$ une famille libre. Il existe une famille
orthonormée $(e_1,\dots,e_p)$ telle que
$\operatorname{Vect}(e_1,\dots,e_k)=\operatorname{Vect}(v_1,\dots,v_k)$
pour tout $k$, construite de proche en proche :
\[
w_k=v_k-\sum_{i<k}\langle v_k,e_i\rangle\,e_i ,
\qquad
e_k=\frac{w_k}{\norme{w_k}} :
\]
on retranche de $v_k$ ses « ombres » sur les directions déjà construites,
puis on norme.
\end{theoreme}

\begin{savoirfaire}[Mener un Gram-Schmidt sans erreur]
Travailler avec les $w_k$ (non normés) le plus longtemps possible, ne
normaliser qu'à la fin ; vérifier chaque orthogonalité au fur et à mesure
— un produit scalaire nul se teste en une ligne. Sur les polynômes avec
$\langle P,Q\rangle=\int_0^1PQ$, l'algorithme appliqué à $(1,X,X^{2})$
produit $1$, $X-\frac12$, $X^{2}-X+\frac16$ : des \emph{polynômes
orthogonaux} (exercices 11, 20, 24) — on y reconnaît le polynôme de
l'exercice 22 du chapitre 15.
\end{savoirfaire}

\essSection{5}{Bases orthonormées}

\begin{theoreme}[Existence et usage]
Tout espace euclidien possède des bases orthonormées, et toute famille
orthonormée se complète en une BON (base orthonormée incomplète —
appliquer Gram-Schmidt à une complétion quelconque). Dans une BON
$(e_1,\dots,e_n)$, tout se calcule par coordonnées :
\[
x=\sum_{i=1}^{n}\langle x,e_i\rangle\,e_i ,
\qquad
\langle x,y\rangle=\sum_{i=1}^{n}\langle x,e_i\rangle\langle y,e_i\rangle ,
\qquad
\norme{x}^{2}=\sum_{i=1}^{n}\langle x,e_i\rangle^{2} :
\]
en BON, tout espace euclidien ressemble à $\R^{n}$ canonique.
\end{theoreme}

\essSection{6}{Projection orthogonale et distance}

\noindent
\begin{minipage}[c]{0.55\linewidth}
\begin{theoreme}[Supplémentaire orthogonal ; projection]
Soit $F$ un sous-espace \emph{de dimension finie} d'un préhilbertien $E$.
Alors $E=F\oplus F^{\perp}$, et la \textbf{projection orthogonale}
$p_F$ sur $F$ s'exprime dans toute BON $(e_1,\dots,e_p)$ de $F$ :
\[
p_F(x)=\sum_{i=1}^{p}\langle x,e_i\rangle\,e_i .
\]
En dimension finie, $\dim F^{\perp}=\dim E-\dim F$ ; un hyperplan est
l'orthogonal d'une droite $\operatorname{Vect}(u)$ : $u$ en est un
\emph{vecteur normal}.
\end{theoreme}
\end{minipage}\hfill
\begin{minipage}[c]{0.41\linewidth}
\centering
\begin{tikzpicture}[scale=0.94,line cap=round]
  \fill[otNavy,opacity=0.06] (-1.8,-0.5) -- (1.3,-0.9) -- (2.0,0.5) -- (-1.1,0.9) -- cycle;
  \draw[otNavy,line width=0.7pt] (-1.8,-0.5) -- (1.3,-0.9) -- (2.0,0.5) -- (-1.1,0.9) -- cycle;
  \fill[otInk] (0.5,2.1) circle (1.6pt);
  \fill[otNavy] (0.3,-0.15) circle (1.6pt);
  \draw[otGreenDark,line width=1pt,dash pattern=on 2.5pt off 1.8pt] (0.5,2.1) -- (0.3,-0.15);
  \draw[otMuted,line width=0.7pt] (0.5,2.1) -- (1.6,-0.35);
  \fill[otMuted] (1.6,-0.35) circle (1.3pt);
  \node[font=\scriptsize,text=otInk,anchor=west] at (0.58,2.15) {$x$};
  \node[font=\scriptsize,text=otNavy,anchor=north] at (0.3,-0.28) {$p_F(x)$};
  \node[font=\scriptsize,text=otGreenDark,anchor=west] at (0.52,1.0) {$d(x,F)$};
  \node[font=\scriptsize,text=otMuted,anchor=north west] at (1.62,-0.38) {$y$};
  \node[font=\scriptsize,text=otNavy] at (-1.42,-0.68) {$F$};
\end{tikzpicture}\\[0.5mm]
\legende{Le projeté $p_F(x)$ est le point de $F$ le plus proche de $x$ :
tout autre $y\in F$ est plus loin (Pythagore).}
\end{minipage}

\begin{theoreme}[Distance à un sous-espace]
$d(x,F)=\norme{x-p_F(x)}$ est la plus petite distance de $x$ à un point de
$F$, atteinte \emph{uniquement} en $p_F(x)$ : pour $y\in F$,
$\norme{x-y}^{2}=\norme{x-p_F(x)}^{2}+\norme{p_F(x)-y}^{2}$ (Pythagore).
Et $\norme{x}^{2}=\norme{p_F(x)}^{2}+d(x,F)^{2}$. Cas de l'hyperplan
$H=\operatorname{Vect}(u)^{\perp}$ :
\[
p_H(x)=x-\frac{\langle x,u\rangle}{\norme{u}^{2}}\,u ,
\qquad
d(x,H)=\frac{|\langle x,u\rangle|}{\norme{u}} .
\]
Projeter, c'est \emph{approcher au mieux} : moindre carrés, polynôme le
plus proche, valeur moyenne — tout le chapitre d'exercices.
\end{theoreme}

%% =====================================================================
%%  EXERCICES
%% =====================================================================
\exercices
\rubrique{Applications}

\begin{exercice}[1]{Produit scalaire ou pas ?}
Les applications suivantes sont-elles des produits scalaires ?
\begin{questions}
  \item Sur $\R^{2}$ :
    $\varphi\bigl((x,y),(x',y')\bigr)=xx'-yy'$, puis
    $\psi\bigl((x,y),(x',y')\bigr)=2xx'+3yy'$.
  \item Sur $\mathcal{C}\bigl(\intFF{0}{1},\R\bigr)$ :
    $\langle f,g\rangle=\displaystyle\int_{0}^{1}fg$, puis
    $\langle f,g\rangle=f(0)\,g(0)$.
\end{questions}
\end{exercice}

\begin{exercice}[1]{Calculs dans $\R^{3}$}
Soient $u=(1,2,2)$ et $v=(2,-1,2)$.
\begin{questions}
  \item Calculer $\norme{u}$, $\norme{v}$, $\langle u,v\rangle$ et le
    cosinus de l'angle entre $u$ et $v$.
  \item Calculer la distance $\norme{u-v}$.
\end{questions}
\end{exercice}

\begin{exercice}[1]{Cauchy-Schwarz en pratique}
\begin{questions}
  \item Montrer que pour tous réels $a_1,\dots,a_n$ :
    $\Bigl(\displaystyle\sum_{i=1}^{n}a_i\Bigr)^{2}
    \leqslant n\displaystyle\sum_{i=1}^{n}a_i^{2}$.
    Quand y a-t-il égalité ?
  \item Vérifier l'inégalité sur $(1,2,3,4)$, et le cas d'égalité sur
    $(2,2,2)$.
\end{questions}
\end{exercice}

\begin{exercice}[1]{Polarisation et parallélogramme}
Soient $u=(2,1,-1)$ et $v=(1,3,2)$.
\begin{questions}
  \item Vérifier l'identité
    $\norme{u+v}^{2}=\norme{u}^{2}+\norme{v}^{2}+2\langle u,v\rangle$.
  \item Retrouver $\langle u,v\rangle$ par polarisation, et vérifier
    l'identité du parallélogramme.
\end{questions}
\end{exercice}

\begin{exercice}[1]{Orthogonalité et Pythagore}
\begin{questions}
  \item Vérifier que $u=(1,2)$ et $v=(-2,1)$ sont orthogonaux, et
    contrôler Pythagore.
  \item Montrer qu'une famille orthogonale de vecteurs non nuls est
    libre.
\end{questions}
\end{exercice}

\begin{exercice}[1]{Gram-Schmidt dans $\R^{3}$}
Orthonormaliser la famille $\bigl((1,1,0),(1,0,1),(0,1,1)\bigr)$
(chapitre 12). On donnera les vecteurs intermédiaires $w_k$ avant
normalisation.
\end{exercice}

\begin{exercice}[1]{Coordonnées en base orthonormée}
Soit la BON de $\R^{2}$ :
$e_1=\bigl(\tfrac35,\tfrac45\bigr)$, $e_2=\bigl(-\tfrac45,\tfrac35\bigr)$.
\begin{questions}
  \item Calculer les coordonnées de $x=(1,2)$ dans cette base.
  \item Vérifier que la somme des carrés des coordonnées vaut
    $\norme{x}^{2}$.
\end{questions}
\end{exercice}

\begin{exercice}[1]{Projection sur une droite}
Soit $D=\operatorname{Vect}(u)$ avec $u=(1,2,2)$, et $x=(3,0,0)$.
\begin{questions}
  \item Calculer le projeté orthogonal de $x$ sur $D$.
  \item En déduire $d(x,D)$, et vérifier
    $\norme{x}^{2}=\norme{p(x)}^{2}+d(x,D)^{2}$.
\end{questions}
\end{exercice}

\rubrique{Entraînement}

\begin{exercice}[2]{Le produit scalaire des matrices}
Sur $\mathcal{M}_2(\R)$, on pose
$\langle A,B\rangle=\operatorname{tr}\bigl(A^{\mathsf T}B\bigr)$.
\begin{questions}
  \item Vérifier que c'est le produit scalaire « coefficient par
    coefficient », et calculer $\langle A,B\rangle$ pour
    $A=\begin{pmatrix}1&2\\3&-1\end{pmatrix}$,
    $B=\begin{pmatrix}0&1\\1&2\end{pmatrix}$.
  \item Montrer que les matrices élémentaires $E_{ij}$ forment une BON.
  \item Quel était le trinôme du chapitre 13, exercice 22, dans ce
    langage ?
\end{questions}
\end{exercice}

\begin{exercice}[2]{Le produit scalaire intégral}
Sur $\mathcal{C}\bigl(\intFF{0}{1},\R\bigr)$, on pose
$\langle f,g\rangle=\displaystyle\int_{0}^{1}fg$.
\begin{questions}
  \item Pourquoi est-ce bien un produit scalaire ? (Chapitre 15.)
  \item Calculer $\norme{x}$, $\langle 1,x\rangle$, et le cosinus de
    l'angle entre les fonctions $1$ et $x$.
\end{questions}
\end{exercice}

\begin{exercice}[2]{Gram-Schmidt sur les polynômes}
Sur $\R_2[X]$ muni de
$\langle P,Q\rangle=\displaystyle\int_{0}^{1}PQ$, orthogonaliser la
famille $(1,X,X^{2})$ (sans normaliser).
\begin{questions}
  \item Montrer qu'on obtient $P_0=1$, $P_1=X-\dfrac12$ et
    $P_2=X^{2}-X+\dfrac16$.
  \item Calculer $\norme{P_1}^{2}$ et $\norme{P_2}^{2}$, et reconnaître
    dans $6P_2$ le polynôme de l'exercice 22 du chapitre 15.
\end{questions}
\end{exercice}

\begin{exercice}[2]{L'orthogonal d'un sous-espace}
Dans $\R^{3}$ canonique, soit $F=\operatorname{Vect}\bigl((1,1,1)\bigr)$.
\begin{questions}
  \item Déterminer $F^{\perp}$, sa dimension et une base.
  \item Vérifier $\R^{3}=F\oplus F^{\perp}$ et
    $\bigl(F^{\perp}\bigr)^{\perp}=F$.
\end{questions}
\end{exercice}

\begin{exercice}[2]{Projection sur un plan}
Soit $H=\ens{(x,y,z)\mid x+y+z=0}$ et $a=(1,2,3)$.
\begin{questions}
  \item Donner un vecteur normal de $H$ et calculer $p_H(a)$.
  \item Calculer $d(a,H)$. Comparer avec la décomposition
    $(1,2,3)=(-1,0,1)+(2,2,2)$ du chapitre 12 : que voit-on de plus ici ?
\end{questions}
\end{exercice}

\begin{exercice}[2]{Distance à un hyperplan}
\begin{questions}
  \item Établir la formule
    $d(x,H)=\dfrac{|\langle x,u\rangle|}{\norme{u}}$ pour
    $H=\operatorname{Vect}(u)^{\perp}$.
  \item Calculer la distance de $(3,1,1)$ au plan $x+2y-2z=0$.
\end{questions}
\end{exercice}

\begin{exercice}[2]{Le projeté en base orthonormée}
Soit $F=\operatorname{Vect}\bigl((1,1,0),(1,0,1)\bigr)$ dans $\R^{3}$ et
$x=(3,0,0)$.
\begin{questions}
  \item Construire une BON de $F$ (exercice 6) et calculer $p_F(x)$.
  \item En déduire $d(x,F)$, et vérifier l'orthogonalité de $x-p_F(x)$
    avec $F$.
\end{questions}
\end{exercice}

\begin{exercice}[2]{La constante la plus proche}
Soit $f$ continue sur $\intFF{0}{1}$. On cherche la constante $c$
minimisant $\displaystyle\int_{0}^{1}\bigl(f(t)-c\bigr)^{2}\mathrm{d}t$.
\begin{questions}
  \item Interpréter le problème comme une projection orthogonale, et
    montrer que la solution est la \emph{valeur moyenne} de $f$
    (chapitre 15).
  \item Expliciter le minimum pour $f(t)=t^{2}$.
\end{questions}
\end{exercice}

\begin{exercice}[2]{Moindres carrés polynomiaux}
Déterminer $a,b$ minimisant
$\displaystyle\int_{0}^{1}\bigl(t^{2}-a-bt\bigr)^{2}\mathrm{d}t$, et la
valeur du minimum. \emph{Projeter $X^{2}$ sur $\R_1[X]$ : les conditions
d'orthogonalité donnent deux équations linéaires.}
\end{exercice}

\begin{exercice}[2]{Pythagore généralisé et inégalité de Bessel}
\begin{questions}
  \item Montrer que pour une famille orthogonale :
    $\norme{x_1+\dots+x_p}^{2}=\sum\norme{x_i}^{2}$.
  \item Soit $(e_1,\dots,e_p)$ orthonormée (pas forcément une base).
    Montrer l'inégalité de Bessel :
    $\displaystyle\sum_{i=1}^{p}\langle x,e_i\rangle^{2}
    \leqslant\norme{x}^{2}$, avec égalité si et seulement si
    $x\in\operatorname{Vect}(e_1,\dots,e_p)$.
\end{questions}
\end{exercice}

\rubrique{Approfondissement}

\begin{exercice}[3]{La meilleure approximation cubique}
Sur $\intFF{0}{1}$, déterminer le polynôme $P\in\R_2[X]$ minimisant
$\displaystyle\int_{0}^{1}\bigl(t^{3}-P(t)\bigr)^{2}\mathrm{d}t$, et
montrer que le minimum vaut $\dfrac{1}{2800}$.
\emph{Trois conditions d'orthogonalité, ou la famille orthogonale de
l'exercice 11.}
\end{exercice}

\begin{exercice}[3]{Les polynômes de Legendre}
Sur $\intFF{-1}{1}$, on pose
$\langle P,Q\rangle=\displaystyle\int_{-1}^{1}PQ$.
\begin{questions}
  \item Vérifier que $1$, $X$, $X^{2}-\dfrac13$, $X^{3}-\dfrac35X$ sont
    deux à deux orthogonaux.
  \item Vérifier sur $n=2$ et $n=3$ la formule de Rodrigues : la dérivée
    $n$-ième de $(X^{2}-1)^{n}$ est proportionnelle au $n$-ième polynôme
    de la famille.
  \item Pourquoi retrouve-t-on, à normalisation près, le Gram-Schmidt de
    $(1,X,X^{2},X^{3})$ ?
\end{questions}
\end{exercice}

\begin{exercice}[3]{Le déterminant de Gram}
Pour $u,v\in E$, on pose
$G(u,v)=\begin{vmatrix}\langle u,u\rangle&\langle u,v\rangle\\
\langle v,u\rangle&\langle v,v\rangle\end{vmatrix}$.
\begin{questions}
  \item Montrer que $G(u,v)\geqslant0$, avec nullité si et seulement si
    $(u,v)$ est liée. Quel théorème retrouve-t-on ?
  \item Dans $\R^{2}$, calculer $G(u,v)$ pour $u=(3,1)$, $v=(1,2)$ et
    comparer à l'aire du parallélogramme (chapitre 14) : montrer qu'en
    général $G(u,v)$ est le carré de l'aire.
\end{questions}
\end{exercice}

\begin{exercice}[3]{Familles obtusangles}
On dit qu'une famille est \emph{obtusangle} si
$\langle x_i,x_j\rangle<0$ pour tous $i\neq j$.
\begin{questions}
  \item Construire trois vecteurs obtusangles dans le plan.
  \item Montrer que si $(x_1,\dots,x_p)$ est obtusangle, alors
    $(x_1,\dots,x_{p-1})$ est libre. \emph{Dans une relation
    $\sum\lambda_ix_i=0$, séparer les coefficients positifs et négatifs
    et prendre un produit scalaire astucieux.}
  \item En déduire qu'un espace euclidien de dimension $n$ contient au
    plus $n+1$ vecteurs obtusangles.
\end{questions}
\end{exercice}

\begin{exercice}[3]{L'inégalité de Hadamard}
Soit $A\in\mathcal{M}_n(\R)$ de colonnes $C_1,\dots,C_n$.
\begin{questions}
  \item Montrer que
    $|\det A|\leqslant\norme{C_1}\cdots\norme{C_n}$.
    \emph{Orthonormaliser les colonnes par Gram-Schmidt et suivre le
    déterminant.}
  \item Quand y a-t-il égalité ? Vérifier sur une matrice diagonale et
    sur $\begin{pmatrix}1&2&0\\0&1&1\\1&0&1\end{pmatrix}$.
  \item Interpréter géométriquement avec les volumes (chapitre 14).
\end{questions}
\end{exercice}

\begin{exercice}[3]{Les racines des polynômes orthogonaux}
Soit $(P_0,P_1,P_2,\dots)$ la famille orthogonale de l'exercice 11
($\deg P_n=n$), pour $\langle P,Q\rangle=\int_0^1PQ$.
\begin{questions}
  \item Montrer que $P_n$ ($n\geqslant1$) est orthogonal à tout polynôme
    de degré $<n$.
  \item En utilisant l'exercice 22 du chapitre 15, montrer que $P_n$
    possède au moins $n$ changements de signe dans $\intOO{0}{1}$ — donc
    exactement $n$ racines réelles \emph{distinctes}, toutes dans
    $\intOO{0}{1}$.
  \item Vérifier pour $P_2$ et pour
    $P_3=X^{3}-\tfrac32X^{2}+\tfrac35X-\tfrac{1}{20}$ (exercice 19).
\end{questions}
\end{exercice}

% !TEX root = ../../main.tex
%% Défi du chapitre 18 — Séries de Fourier et inégalité de Bessel
%% Résultats vérifiés par verifications/defis_16_20.py.

\rubrique{Défi}

\begin{defi}{Séries de Fourier et inégalité de Bessel}{2 h}
On munit $E=\mathcal{C}\bigl(\intFF{-\pi}{\pi},\R\bigr)$ de
$\langle f,g\rangle=\dfrac1\pi\displaystyle\int_{-\pi}^{\pi}f(x)g(x)\,\mathrm{d}x$.
Approcher une fonction par des sommes de sinus et cosinus — l'idée de
Joseph Fourier (1807) — revient à la projeter orthogonalement sur des
espaces de \emph{polynômes trigonométriques}.

\partieExo{A}{Une famille orthonormée}
\begin{questions}
  \item Montrer que $\langle\cdot,\cdot\rangle$ est un produit scalaire sur
    $E$.
  \item Pour $k\geqslant1$, on note $c_k:x\mapsto\cos kx$ et
    $s_k:x\mapsto\sin kx$, et $c_0$ la fonction constante
    $\frac{1}{\sqrt2}$. Montrer, en linéarisant les produits, que la
    famille $(c_0,c_1,s_1,\dots,c_n,s_n)$ est orthonormée.
  \item On note $T_n$ l'espace qu'elle engendre. Montrer que la projection
    orthogonale de $f$ sur $T_n$ est
    \[
    S_n(f)=\frac{a_0}{2}+\sum_{k=1}^{n}\bigl(a_k\cos kx+b_k\sin kx\bigr),
    \]
    où
    $a_k=\frac1\pi\int_{-\pi}^{\pi}f(x)\cos kx\,\mathrm{d}x$ et
    $b_k=\frac1\pi\int_{-\pi}^{\pi}f(x)\sin kx\,\mathrm{d}x$.
\end{questions}

\partieExo{B}{Le cas de $f(x)=x$}
\begin{questions}[resume]
  \item Montrer que $a_k=0$ pour tout $k$ et que
    $b_k=\dfrac{2(-1)^{k+1}}{k}$.
  \item Calculer $\norme{f}^{2}$, puis déduire de l'inégalité de Bessel
    (chapitre 18) que
    $\displaystyle\sum_{k=1}^{n}\frac{1}{k^{2}}\leqslant\frac{\pi^{2}}{6}$
    pour tout $n$. Qu'en déduit-on pour
    $\zeta(2)=\sum_{k\geqslant1}\frac{1}{k^{2}}$ ?
  \item Montrer que
    $d(f,T_n)^{2}=\dfrac{2\pi^{2}}{3}-4\displaystyle\sum_{k=1}^{n}\frac{1}{k^{2}}$.
    On admet que $d(f,T_n)\to0$ (égalité de Parseval, programme de deuxième
    année) : en déduire $\zeta(2)=\dfrac{\pi^{2}}{6}$. (Le défi du
    chapitre 19 en donne une preuve complète.)
\end{questions}

\medskip
\begin{center}
\begin{tikzpicture}[x=1.05cm,y=0.62cm,line cap=round]
  \draw[otRule,line width=0.3pt] (-3.1416,-3.1416) -- (-3.1416,3.1416) (3.1416,-3.1416) -- (3.1416,3.1416);
  \draw[otMuted,line width=0.5pt,-{Stealth[length=1.6mm]}] (-3.5,0) -- (3.6,0);
  \draw[otMuted,line width=0.5pt,-{Stealth[length=1.6mm]}] (0,-4) -- (0,4.1);
  \draw[otInk,line width=1.4pt] (-3.1416,-3.1416) -- (3.1416,3.1416);
  \draw[otNavy!45,line width=0.8pt] plot[domain=-3.1416:3.1416,samples=90] ({\x},{2*sin(deg(\x))});
  \draw[otNavy!75,line width=0.9pt] plot[domain=-3.1416:3.1416,samples=150]
    ({\x},{2*(sin(deg(\x))-sin(deg(2*\x))/2+sin(deg(3*\x))/3)});
  \draw[otGreenDark,line width=1.1pt] plot[domain=-3.1416:3.1416,samples=300]
    ({\x},{2*(sin(deg(\x))-sin(deg(2*\x))/2+sin(deg(3*\x))/3-sin(deg(4*\x))/4
      +sin(deg(5*\x))/5-sin(deg(6*\x))/6+sin(deg(7*\x))/7)});
  \node[font=\tiny,text=otMuted,below] at (3.1416,-0.05) {$\pi$};
  \node[font=\tiny,text=otMuted,below] at (-3.1416,-0.05) {$-\pi$};
  \node[font=\scriptsize,text=otInk,anchor=west] at (3.2,3.2) {$f(x)=x$};
  \node[font=\scriptsize,text=otNavy!60,anchor=south] at (1.57,2.05) {$S_1$};
  \node[font=\scriptsize,text=otNavy!80,anchor=north] at (2.3,1.2) {$S_3$};
  \node[font=\scriptsize,text=otGreenDark,anchor=west] at (2.85,3.7) {$S_7$};
\end{tikzpicture}\\[0.5mm]
\legende{Les projections $S_1$, $S_3$, $S_7$ de $f(x)=x$ : elles épousent
$f$ de mieux en mieux à l'intérieur, mais retombent à $0$ en $\pm\pi$, avec
un dépassement persistant près des bords.}
\end{center}

\partieExo{C}{Convergence ponctuelle et phénomène de Gibbs}
\begin{questions}[resume]
  \item Montrer que
    $S_n(f)\bigl(\frac\pi2\bigr)=2\sum_{j}\frac{(-1)^{j}}{2j+1}$, la somme
    portant sur les $j$ tels que $2j+1\leqslant n$. En admettant que
    $S_n(f)\bigl(\frac\pi2\bigr)\to f\bigl(\frac\pi2\bigr)$, retrouver la
    formule de Leibniz
    $1-\frac13+\frac15-\frac17+\cdots=\frac\pi4$ (chapitre 19).
  \item Calculer $S_n(f)(\pi)$ pour tout $n$ et commenter la figure : le
    maximum de $S_n(f)$ près de $\pi$ ne tend pas vers $\pi$ mais vers
    environ $3{,}70$ (on l'admet). Pourquoi cela n'est-il pas en
    contradiction avec $d(f,T_n)\to0$ ?
\end{questions}
\end{defi}

